\section{Related Works}
\label{section:3_relatedWorks}
\subsection{Text-Guided Image Editing}
\noindent\textbf{Training-based methods.}
Training-based approaches learn editing capabilities from paired or synthetically constructed datasets. InstructPix2Pix~\cite{InstructPix2Pix} fine-tunes a conditional diffusion model on synthetic instruction--image pairs, enabling single-forward-pass editing from natural language instructions. Imagic~\cite{Imagic} performs complex non-rigid edits by optimizing the text embedding and fine-tuning the model, but requires expensive per-image optimization. BrushNet~\cite{BrushNet} and BrushEdit~\cite{BrushEdit} address region-based inpainting through dual-branch architectures and automated mask generation, respectively. FLUX.1 Kontext~\cite{FLUX.1Kontext} unifies generation and editing by concatenating input and output tokens in latent space with a flow matching objective. While these methods achieve strong editing quality, they face the additional challenge of training overhead and data dependency.
\\
\noindent\textbf{Inversion-based training-free methods.}
A major line of training-free editing relies on inverting a real image into the generative model's latent space. Prompt-to-Prompt (\textit{P2P})~\cite{prompt2prompt} demonstrates that cross-attention maps in diffusion models govern the spatial relationship between image layout and text tokens, enabling localized editing by injecting source attention maps into the target generation process; however, \textit{P2P} operates solely in a text-to-image setting and cannot directly edit real images. Null-text Inversion~\cite{Null-text} addresses this limitation by reconstructing the noise trajectory of an input image through DDIM inversion and optimizing the null-text embedding used in classifier-free guidance. Prompt Tuning Inversion~\cite{Prompt-Tuning} encodes input image information into a learnable conditional embedding via prompt tuning, then interpolates between the optimized and target embeddings to balance editability and fidelity. PnP Inversion~\cite{PnP_Inversion} further simplifies inversion-based editing by rectifying inversion deviations directly within the source diffusion branch while leaving the target branch unaltered. LEDits++~\cite{LEDits++} proposes an efficient perfect inversion scheme with implicit masking that limits changes to semantically relevant regions, supporting multiple simultaneous edits. RF-Inversion~\cite{RF-inversion} extends inversion to rectified flow models by formulating it as dynamic optimal control via a linear quadratic regulator, enabling zero-shot editing for models such as FLUX. Despite progressive improvements, all inversion-based methods remain fundamentally constrained by the accuracy of the inverted trajectory.

\noindent\textbf{Non-inversion training-free methods.}
An alternative line of work avoids explicit inversion altogether. SDEdit~\cite{SDEdit} adds noise to the input image and iteratively denoises through a stochastic differential equation prior, naturally balancing realism and faithfulness but offering limited spatial control. AREdit~\cite{AREdit} is the first to apply training-free editing to Visual Autoregressive Models, caching token indices and probability distributions from the source generation to bridge the information gap between text prompts and the input image. Our method falls in this non-inversion category but differs from AREdit in its editing mechanism: whereas AREdit replays cached token statistics (requiring per-category hyperparameter tuning), we use cross-modal attention maps to explicitly construct spatial editing masks, providing an adaptive and tuning-free alternative.
\subsection{Visual Autoregressive Modeling}
Early discrete visual representation learning established the foundation for autoregressive image generation. VQ-VAE~\cite{VQVAE} introduces vector quantization into the VAE bottleneck, learning discrete latent codes via a straight-through estimator and modeling their distribution with an autoregressive prior; its discrete spatial codes are a conceptual precursor to the token-level editing operations our method performs. VQ-GAN~\cite{VQGAN} extends this by incorporating adversarial training and perceptual losses, producing a perceptually faithful codebook that enables high-resolution image synthesis when combined with an autoregressive transformer over the flattened code grid.
A paradigm shift occurs with VAR~\cite{VAR}, which redefines autoregressive image generation as coarse-to-fine next-scale prediction rather than raster-scan next-token prediction. A multi-scale VQVAE tokenizer encodes images into hierarchical token maps at increasing resolutions, and a transformer autoregressively predicts all tokens at the next finer scale conditioned on coarser scales. This formulation yields substantial improvements in both quality and speed, and demonstrates clear power-law scaling behavior analogous to large language models. Critically for our purposes, the spatial locality of each token---each corresponding to a well-defined patch at a specific scale---is what makes selective token replacement a geometrically meaningful editing operation. HART~\cite{HART} further introduces a hybrid tokenizer that decomposes visual representations into discrete tokens for coarse structure and continuous residual tokens for fine details, with the latter modeled by a lightweight residual diffusion module, bridging the quality gap between discrete autoregressive and continuous diffusion approaches. Infinity~\cite{Infinity} extends VAR by replacing finite-vocabulary codebooks with bitwise token prediction using an infinite-vocabulary tokenizer paired with a bitwise self-correction mechanism, enabling stronger scaling behavior and high-resolution generation. Complementarily, BSQ-VAE~\cite{BSQVAE} proposes Binary Spherical Quantization, a codebook-free approach that projects visual embeddings onto a hypersphere and applies binary quantization, achieving extreme compression with minimal distortion; its lossless encoding property is directly exploited in our framework, where transplanting BSQ-VAE tokens from a source image guarantees pixel-faithful background preservation without any reconstruction distortion.
Our work builds upon the VAR architecture, exploiting two of its structural properties that continuous generative paradigms do not share: the spatial locality of discrete tokens, which enables geometrically precise editing masks, and the self-consistent invertibility of BSQ-VAE encoding, which enables background preservation by construction rather than by approximation.